\phantomsection\addcontentsline{toc}{section}{Introducción a los Transformers y su relación 
con los LLM}
\section*{Introducción a los Transformers y su relación 
con los LLM}

Los modleos de lenguaje han recorrido un largo camino desde los enfoques estadísticos basados
en conteos de n-gramas hasta convertirse en los actuales modelos de lenguaje a gran escala, 
también conocidos como LLMs por sus siglas en inglés. En ese sentido, la arquitectura
Transformer (Vaswani et al., 2017) ha sido el principal punto de partida, con el cual se revolucionó la manera en
que presentamos y procesmaos secuencias linguísticas. 

A diferencia de los modelos recurrentes, como las RNN, LSTM o GRU, que procesan el texto paso 
a paso, y que sufrían de limitaciones para capturar dependencias largas, los Transformers 
incorporaron un mecanismo esencial para su éxito: la auto-atención (o \textit{Self-Attention}).
Este mecabnismo permite que cada palabra o token en una secuencia "se fije" de manera
directa en cualquier otro, sin importar la distancia que los separe, construyendo representaciones
contextuales mucho más ricas y globales. 

Este avance fue toda una revolución para el Procesamiento de Lenguaje Natural, pues creó 
modelos que podían manejar contextos extensos y aprender patrones complejos de manera eficiente. 
Como resultado, los Transformers se convirtieron muy rápidamente en la arquitectura dominante
en prácticamente todas las tareas del NLP, desde la traducción automática, hasta la generación 
de texto. 

Los LLMs contemporáneos como GPT, LLaMA, Gemini o DeepSeek no son más que instancias a gran
escala de esta arquitectura. Es decir, lo que se ha hecho es tomar la arquitectura y volverla
gigantesca, de tal modo que ahora tenemos modelos sumamente eficaces. La característrica de
\textit{Large} de los LLMs se refiere tanto a su número de parámetros como a la magnitud
del corpus con el que fue entrenado. Este diseño, que no necesita de convoluciones ni de 
recurrencia, representa un cambio de paradigma en la modelación del lenguaje. Se trata 
de una arquitectura altamente paralelizable y que puede capturar dependencias globales en el
texto de manera directa. 

\subsection{El cambio de paradigma}

El Transformer (Vaswani et al., 2017), es una arquitectura de red neuronal diseñada originalmente
para tareas de traducción automática, pero que pronto se convirtió en el estándar para el NLP, 
encontrando aplicaciones en muchos campos. 

A diferencia de los modelos recurrentes, que procesan de manera secuencial el texto, el Transformer
elimina la recurrencia y se basa en bloques paralelos compuestos de:

\begin{itemize}
    \item Mecanismo de self-atention o multi-head attention.
    \item Redes feed-forward posicionales.
    \item Normalización por capas y conexiones residuales.
\end{itemize}

En general, un Transformer puede adaptar tres variantes arquitectónicas principales:

\begin{itemize}
    \item \textbf{Decoder-only}: utilizado en modelos generativos como GPT, LLaMA o Mistral.
    \item \textbf{Encoder-only}: especializado en clasificación y comprensión como BERT, BETO o RoBERTa. 
    \item \textbf{Encoder-decoder}: comúnmente utilizados para traducción y tareas secuencia a secuencia.
\end{itemize}

Una de sus grandes ventajas es que pueden procesar contexto completo en paralelo, y aprender
dependencias de largo alcance en un solo paso, como no podrían los modelos recurrentes. 

\subsection{\textit{Self-Attention}}

Dado un conjunto de representaciones entrada $(x_1, x_2, \dots, x_n) \in \mathbb{R}^d$, una por 
cada token, la auto-atención produce nuevas representaciones $a_1, a_2, \dots, a_n$ donde cada
$a_i$ es una combinación ponderada de todos los vectores $x_j$ donde $j < n$.

La idea es que un token puede consultar información de otros tokens en función de su similitud
semántica y contextual.

Se definen tres proyecciones aprendidas de cada vector de entrada:

\[
    q_i = x_i W^Q, \;\;\;\;\;\;\; k_i = x_i W^K, \;\;\;\;\;\;\; v_i = x_i W^V
\]

Donde $W^Q, W^K, W^V \in \mathbb{R}^{d \times d_k}$ son matrices de parámetros.

La similaridad entre token $i$ (consulta) y otro token $j$ (clave) se calcula mediante un producto
punto escalado:

\[
    \text{score}(i,j) = \frac{q_i \cdot k_j}{\sqrt{d_k}}
\]

El uso del factor $\sqrt{d_k}$ estabiliza los gradientes y evita valores extremos al aplicar
la exponenciación en el softmax. Estas puntuaciones se normalizan con una Softmax, obteniendo
los pesos de atención:

\[
    \alpha_{ij} = \frac{\exp \Big(\text{score}(i,j)\Big)}{\sum_{m=1}^{n} \exp \Big(\text{score}(i,m)\Big)}
\]

Finalmente, la represebtación contextual del token $i$ se obtiene como una suma ponderada
de los valores:

\[
    \alpha_i = \sum_{j=1}^{n} \alpha_{ij} v_j
\]

Este mecanismo se aplica en paralelo para todos los tokens de la secuencia, generando la matriz
de salidas $A \in \mathbb{R}^{n \times d_v}$.

\subsection{\textit{Multi-head Attention}}

En la pr\'actica, no se utiliza una sola cabeza de atenci\'on, sino m\'ultiples. Con $h$ cabezas, se calculan diferentes subespacios de atenci\'on:

$$\text{head}_h = \text{Attention}(XW^Q_h, XW^K_h, XW^V_h)$$

y luego se concatenan todas las salidas para formar la representaci\'on final:

$$\text{MHA}(X) = [\text{head}_1; \dots; \text{head}_h] W^O$$

donde $W^O$ es una matriz de proyecci\'on adicional.

\subsection{Transformers populares}

La introducci\'on de la arquitectura Transformer no s\'olo transform\'o la investigaci\'on en PLN, sino que tambi\'en dio origen a una serie de modelos que hoy son considerados referentes en la evoluci\'on de los modelos de lenguaje a gran escala (\textbf{LLMs}). Entre ellos destacan las familias \textbf{GPT}, \textbf{LLaMA} y m\'as recientemente \textbf{DeepSeek}, cada una con enfoques y motivaciones distintas.

\subsubsection{GPT (Generative Pretrained Transformer)}

Los modelos GPT, desarrollados por OpenAI, constituyen el ejemplo paradigm\'atico de los Transformers decoder-only. Su entrenamiento se basa en la tarea de \textit{next-word prediction}, es decir, la predicci\'on autoregresiva de la siguiente palabra dada una secuencia de entrada.

\begin{itemize}
    \item \textbf{GPT-1 (2018)} introdujo la idea de preentrenar un gran modelo en corpus generales y luego ajustarlo (\textit{fine-tuning}) para tareas espec\'ificas.
    \item \textbf{GPT-2 (2019)} escal\'o el tama\~no (1.5B par\'ametros) y mostr\'o por primera vez capacidades generativas sorprendentes, aunque inicialmente fue liberado de forma parcial por preocupaciones de seguridad.
    \item \textbf{GPT-3 (2020)}, con 175B par\'ametros, consolid\'o el paradigma de los LLMs al mostrar \textit{emergent abilities}: razonamiento b\'asico, traducci\'on, generaci\'on de c\'odigo y di\'alogo. Su ajuste posterior dio origen a ChatGPT, que incorpor\'o \textit{instruction tuning} y \textit{reinforcement learning from human feedback} (RLHF) para mejorar su alineaci\'on con usuarios.
\end{itemize}

El impacto de GPT fue doble: estableci\'o la viabilidad de los LLMs en m\'ultiples dominios y mostr\'o que un \'unico modelo, entrenado con suficiente escala, puede generalizar a una amplia variedad de tareas ling\"u\'isticas sin redise\~nar la arquitectura.

\subsubsection{LLaMA (Large Language Model Meta AI)}

En 2023, Meta present\'o la familia LLaMA, concebida como un esfuerzo por abrir a la comunidad modelos competitivos pero con un enfoque en eficiencia y accesibilidad.

\begin{itemize}
    \item \textbf{LLaMA-1} ofrec\'ia variantes de 7B a 65B par\'ametros entrenadas con datos cuidadosamente filtrados.
    \item \textbf{LLaMA-2 (2023)} introdujo mejoras en el preentrenamiento y versiones afinadas para di\'alogo.
    \item \textbf{LLaMA-3 (2024)} alcanz\'o mayores tama\~nos (hasta 405B par\'ametros en versiones privadas), con una fuerte orientaci\'on a la investigaci\'on abierta y la compatibilidad con entornos de \textit{fine-tuning}.
\end{itemize}

Una de las principales contribuciones de la familia LLaMA fue democratizar el acceso a modelos de calidad comparable a GPT-3, pero bajo licencias abiertas, favoreciendo la investigaci\'on acad\'emica y el desarrollo de aplicaciones personalizadas.

\subsubsection{DeepSeek}

En el ecosistema m\'as reciente ha surgido DeepSeek, desarrollado en China como un competidor directo de modelos occidentales de gran escala. Su relevancia radica en dos aspectos:

\begin{itemize}
    \item \textbf{Optimizaci\'on computacional}: uso de t\'ecnicas avanzadas de entrenamiento distribuido que permiten reducir costos energ\'eticos y de hardware.
    \item \textbf{Competencia global}: representa un movimiento estrat\'egico en la carrera internacional de LLMs, mostrando que la arquitectura Transformer puede adaptarse a entornos con recursos limitados pero ambiciones de despliegue masivo.
\end{itemize}

DeepSeek no s\'olo apunta a competir en calidad, sino tambi\'en a abrir el camino a modelos m\'as eficientes, lo que sugiere que el futuro de los LLMs no depende \'unicamente de escalar par\'ametros, sino de optimizar su entrenamiento y uso.

\subsubsection{Otros modelos influyentes}

Aunque GPT, LLaMA y DeepSeek ocupan hoy el centro de atenci\'on, existen otras familias de Transformers fundamentales en la evoluci\'on del campo:

\begin{itemize}
    \item \textbf{BERT (2018)}: modelo \textit{encoder-only} basado en \textit{masked language modeling}, dise\~nado para tareas de comprensi\'on como clasificaci\'on, b\'usqueda sem\'antica y an\'alisis de sentimiento.
    \item \textbf{T5 (2019)}: modelo \textit{encoder-decoder} que formul\'o todas las tareas de PLN como problemas de traducci\'on de texto a texto.
    \item \textbf{Mistral (2023)}: modelos ligeros de arquitectura \textit{decoder-only}, altamente eficientes, dise\~nados para ser competitivos con menos par\'ametros.
    \item \textbf{Gemini (2023, Google DeepMind)}: sucesor de PaLM, con capacidades multimodales (texto e imagen) y mayor integraci\'on con entornos de b\'usqueda.
\end{itemize}